\Needspace{11\baselineskip}
\section{Enclosing area, obstacle minimizers, and metric variation}
\label{rig4:sec:enclosures}
\label{rig4:enc:section}

The cuts in Definition~\ref{four:def:cuts} count the entire intrinsic
boundary, including every portion coincident with the original boundary.
We identify their area infimum with a full-perimeter minimum. The resulting
enclosure will be used in the numerical inequality. Its variation formula,
which retains all tied minimizers, will be used in the equality argument.

Throughout this section, $(X,g)$ is a smooth connected orientable
four-dimensional Riemannian manifold with nonempty compact smooth
boundary $B$. It is complete with $B$ included and has exactly one end,
which is asymptotically Euclidean. In coordinates on that end we assume
$g-\delta=O_2(r^{-q})$ for some $q>1$. No scalar-curvature assumption is
needed in this section. Write
\[
 a_g(B)=\inf_{\Gamma}\Area_g(\Gamma),
\]
where $\Gamma$ ranges over the full smooth enclosing cuts of
Definition~\ref{four:def:cuts}.

\subsection{A perimeter problem that counts the original boundary}
\label{rig4:enc:obstacle-setup}

Choose a smooth manifold $\widetilde X$ without boundary by attaching an
inner side along $B$, and extend $g$ smoothly through $B$. For example,
one may use the topological double, with a smooth metric extension in a
collar and an arbitrary smooth metric farther inside. Let $O$ be the
closed inner side, so that $\partial O=B$ and
$\widetilde X\setminus O=\operatorname{int}X$. Neither completeness nor a
curvature condition on the extension will be used. We continue to write
$g$ for this auxiliary extension.

Let $\mathcal A$ consist of sets $E$ of locally finite perimeter in
$\widetilde X$ such that $O\subset E$ up to a null set and
$E\setminus O$ is contained, up to a null set, in some compact subset
of $X$.
Sets are identified if their indicator functions agree almost
everywhere. The perimeter is always
\begin{equation}
 P_g(E)=|D\mathbf 1_E|_g(\widetilde X),
 \label{rig4:enc:full-perimeter}
\end{equation}
computed in the extended manifold. Thus $P_g(O)=\Area_g(B)$, and
perimeter on $B$ is retained. Although $O$ itself need not have finite
volume, the competitors differ from it only in a compact exterior
region, so the compactness arguments below are local BV arguments on a
fixed bounded region. We use the standard perimeter compactness,
Gauss--Green, and submodularity facts for BV sets
\cite{AmbrosioFuscoPallara2000,Maggi2012}.

If $D$ is a smooth exterior domain from Definition~\ref{four:def:cuts}, its
filled side is
\[
 E_D=O\cup\bigl(X\setminus\operatorname{int}D\bigr).
\]
Here $\operatorname{int}D$ denotes the intrinsic manifold interior, which
lies in $\operatorname{int}X$. The distant end is absent from $E_D$, and
\begin{equation}
 E_D\in\mathcal A,
 \qquad P_g(E_D)=\Area_g(\partial_{\mathrm{man}}D).
 \label{rig4:enc:cut-to-set}
\end{equation}
The identity includes any frontier on $B$. In particular $D=X$
corresponds to $E_D=O$, with perimeter $\Area_g(B)$.

\begin{proposition}[Full-perimeter minimizing enclosures]
\label{rig4:enc:minimizing-hull}
The minimum of $P_g$ over $\mathcal A$ exists and equals $a_g(B)>0$.
Every minimizer has a regular representative whose frontier $T$ is a
nonempty compact embedded $C^{1,1}$ hypersurface, smooth and minimal on
$T\setminus B$. Its graph functions belong locally to $W^{2,\infty}$,
and hence to $W^{2,2}$, including at contact with $B$. Its outward
coorientation points from the filled set toward a connected exterior
containing the asymptotic end. All minimizing frontiers lie in a fixed
compact subset of $X$.

Every such frontier is a limit in $C^1$, and in area, of admissible smooth
enclosing cuts lying in $\operatorname{int}X$. If in addition
$H_B<0$ for the normal into $X$, every minimizer contains a fixed collar
of $B$ on the exterior side. In this case $T$ is itself a smooth
minimal enclosing cut, and its closed exterior is smooth, connected,
complete with its nonempty boundary included, and one-ended. The
boundary of this exterior is outer area-minimizing there, with every
boundary component counted.
\end{proposition}

\begin{proof}
\emph{A common outer truncation.}
On the coordinate end, let
$Y=\nabla r/|\nabla r|_g$ be the unit normal pointing toward increasing
radius. The differentiated asymptotic decay gives
\begin{equation}
 \operatorname{div}_gY=\frac{3}{r}+O(r^{-1-q})>0
 \qquad (r\ge R_0)
 \label{rig4:enc:outer-mean-convexity}
\end{equation}
after increasing $R_0$. For a competitor $E$ and almost every $R>R_0$,
put $E_R=E\setminus\{r>R\}$. Gauss--Green on
$V_R=E\cap\{r>R\}$ gives
\[
 \int_{V_R}\operatorname{div}_gY\,dV_g
 =\int_{\partial^*E\cap\{r>R\}}
       g(Y,\nu_E)\,dA_g
   -\int_{\{r=R\}}\mathbf 1_E\,dA_g.
\]
The trace on the slicing sphere is understood at a good slicing radius;
$\nu_E$ is the measure-theoretic outward normal of the filled set.
Consequently
\begin{equation}
 P_g(E_R)-P_g(E)
 \le -\int_{V_R}\operatorname{div}_gY\,dV_g\le0,
 \label{rig4:enc:outer-truncation}
\end{equation}
and the inequality is strict if $V_R$ has positive volume.

Fix $R_1>R_0$. For each competitor separately, choose a good radius in
$(R_0,R_1)$ and use Equation~\eqref{rig4:enc:outer-truncation}. The infimum
therefore equals the infimum in the class with $E\setminus O$ confined
to the same fixed truncation $\{r\le R_1\}$ together with the compact
core. BV compactness on a slightly larger compact region and lower
semicontinuity give a minimizer. The constraints of containing $O$ and
being empty in the prescribed exterior tail are closed under this
convergence. The competitor $O$ gives a finite upper bound.

The resulting set is also a global minimizer in $\mathcal A$. Applied
to any global minimizer, Equation~\eqref{rig4:enc:outer-truncation} forbids
positive filled volume beyond a good radius above $R_0$. Choosing such
a radius below a fixed number $R_2$ with $R_0<R_2<R_1$ shows that all
minimizers have their filled exterior parts in the same compact
truncation at $R_2$. This also keeps their frontiers away from the
artificial boundary used in the compactness argument. There is no
requirement that one slicing radius be good for every BV set.

\emph{Regularity at the obstacle.}
A global minimizer is locally perimeter-minimizing among competitors
that contain the smooth open obstacle $O^\circ$. The relevant form of
the Riemannian obstacle regularity theorem is the following: in a smooth
Riemannian manifold of dimension less than eight, a local perimeter
minimizer constrained to contain an open set with $C^2$ boundary has a
$C^{1,1}$ frontier and is smooth away from contact. This is precisely
Regularity Theorem~1.3(iii) of Huisken--Ilmanen
\cite[pp.~368--369]{HuiskenIlmanen2001}, applied with zero bulk term.
Its hypotheses are local, so the arbitrary geometry of the attached
inner side and its possible noncompactness are irrelevant. Here the
ambient dimension is four. The result describes the entire frontier
of the regular representative, not just the reduced boundary outside
an unaccounted exceptional set.

For clarity, the obstacle also gives the usual two-sided almost-minimality
estimate needed for density and compactness estimates. Choose a smooth
compactly supported vector field $Z$ on $\widetilde X$ with $|Z|_g\le1$
and $Z=\nu_O$ on $B$. Such a field is obtained from a tubular
neighborhood of the compact smooth boundary. For a local replacement
$F$, with $F\triangle E$ compactly supported, $F\cup O$ is admissible.
Perimeter submodularity and Gauss--Green imply
\begin{align}
 P_g(E)
 &\le P_g(F\cup O)
 \le P_g(F)+P_g(O)-P_g(F\cap O),\nonumber\\
 P_g(O)-P_g(F\cap O)
 &\le \int_{O\setminus F}\operatorname{div}_gZ\,dV_g
 \le C\Vol_g(E\triangle F).
 \label{rig4:enc:almost-minimality}
\end{align}
Thus local replacements satisfy the standard perimeter
almost-minimality inequality. In particular, the usual perimeter and
volume density estimates apply, including at contact
\cite{Maggi2012}. The $C^{1,1}$ conclusion above directly gives the
asserted Sobolev regularity in graph charts. Off the obstacle, arbitrary
local replacements are allowed; first variation gives zero mean
curvature and interior regularity gives smoothness.

Take the representative that is the closure of its measure-theoretic
interior. Its frontier is the compact $C^{1,1}$ hypersurface just
described, and the filled and exterior sides occupy the two local sides
of every frontier chart. This choice discards all irrelevant null-set
slits. The frontier is nonempty: the filled set contains the open
inner side, whereas its complement contains an exterior tail. On the
connected extended manifold a BV indicator with zero perimeter is
constant almost everywhere, which is impossible here. The nonempty
regular frontier consequently has positive area.

\emph{Connectedness of the exterior.}
The complement of the regular filled set has a component containing the
entire sufficiently distant end. Any other component is relatively
compact, because the frontier is compact and $X$ has one end. If such
a component existed, adjoining it to the filled set would remove its
nonempty boundary and introduce no new perimeter. More explicitly,
near each frontier point there is only one exterior side; the boundary
of a complement component is therefore a union of connected components
of the compact embedded frontier. Filling a bounded component removes
a positive amount of full perimeter, including any portion on $B$.
This contradicts minimality. Thus the exterior is connected.

\emph{Equality with the smooth-cut infimum.}
Equation~\eqref{rig4:enc:cut-to-set} gives
$\min_{\mathcal A}P_g\le a_g(B)$. To prove the reverse inequality, let
$T$ be a minimizing frontier with its outward coorientation. Choose a
smooth compactly supported vector field $V$ with
$g(V,\nu_T)\ge c_0>0$ on $T$; approximate its continuous unit normal in
local charts and use a partition of unity. If $\Phi_t$ is its flow,
compactness and transversality give a $C^1$ flow collar
\[
 \Psi:T\times(-\delta,\delta)\longrightarrow\widetilde X,
 \qquad \Psi(p,t)=\Phi_t(p),
\]
in which the filled side is $\{t\le0\}$. For sufficiently small
$\varepsilon>0$ the displaced filled set satisfies
\[
 E_\varepsilon:=\Phi_\varepsilon(E)
   =E\cup\Psi\bigl(T\times(0,\varepsilon]\bigr).
\]
This inclusion is global: a point can change its membership only when
its flow line crosses $T$, and every such short crossing is outward in
the chosen collar. Thus
$O\subset E\subset\operatorname{int}E_\varepsilon$.
The displaced frontier $T_\varepsilon=\partial E_\varepsilon$ is
disjoint from $O$ and lies in $\operatorname{int}X$. Its distance from
the compact original boundary is positive.

Choose a smaller cooriented collar of $T_\varepsilon$ with compact
closure disjoint from $O$, and a $C^1$ defining function there whose
derivative along $V$ is positive. Smooth this defining function in
$C^1$. A sufficiently close approximation still has positive
$V$-derivative, and its zero set meets each collar fiber exactly once.
It is consequently a smooth hypersurface given by a small graph over
$T_\varepsilon$ along those fibers. The graph displacement extends to
a collar isotopy: in collar coordinates use
$(p,t)\mapsto(p,t+\tau\chi(t)w(p))$, where $w$ is the small graph height,
$\chi$ equals one near zero and vanishes near the collar ends, and
$0\le\tau\le1$. Smallness of $w$ makes the derivative in $t$ positive.
This isotopy is the identity on $O$ and in the distant end. It
therefore preserves enclosure and connectedness of the exterior.
Its final smooth frontier bounds a connected smooth exterior domain,
closed in $X$, and is exactly that domain's full intrinsic boundary.

The resulting cuts are admissible. Their areas tend first to the
area of $T_\varepsilon$ by $C^1$ approximation and then to
$\Area_g(T)$ as $\varepsilon\downarrow0$. Hence
$a_g(B)\le P_g(E)$, proving equality of the two minimization problems
and the claimed approximation. This also proves that the minimum is
independent of the chosen inner extension.

\emph{A strict inner barrier.}
Suppose $H_B<0$ for the normal into $X$. Compactness of $B$ gives an
exterior signed-distance collar $0\le t\le\varepsilon$ whose unit
normal $Y=\partial_t$ satisfies
\[
 \operatorname{div}_gY\le-\beta<0
\]
for fixed $\varepsilon,\beta>0$. For almost every
$s\in(0,\varepsilon)$, let
$O_s=O\cup\{0<t<s\}$ and $W_s=O_s\setminus E$. Extend the collar
field smoothly across $B$ when applying Gauss--Green. The boundary
flux identity on $W_s$ gives
\[
 \int_{W_s}\operatorname{div}_gY\,dV_g
 =\int_{\{t=s\}}\mathbf 1_{E^c}\,dA_g
   -\int_{\partial^*E\cap\{0\le t<s\}}
          g(Y,\nu_E)\,dA_g.
\]
The second integral includes the frontier at $t=0$. Therefore
\begin{equation}
 P_g(E\cup O_s)-P_g(E)
 \le\int_{W_s}\operatorname{div}_gY\,dV_g
 \le-\beta\Vol_g(W_s).
 \label{rig4:enc:inner-truncation}
\end{equation}
Minimality forces $\Vol_g(W_s)=0$. For each minimizer choose a good
$s\in(\varepsilon/2,\varepsilon)$. Its regular representative then
contains the whole open collar $0\le t<\varepsilon/2$ as interior
points in $\widetilde X$. Thus the same collar is included by all
minimizers, despite allowing their good slicing levels to differ.

The frontier is now disjoint from $B$, hence is smooth and minimal.
Its closed exterior $D$ is a smooth submanifold with boundary $T$, closed
as a subset of $X$, with connected interior and the same distant end.
It is complete in its intrinsic length metric. Indeed an intrinsic
Cauchy sequence is Cauchy for the original metric distance and converges
to a point of the closed set $D$ by completeness of $X$. Interior
charts and smooth boundary half-ball charts identify the two local
notions of convergence, with locally comparable path lengths, so the
sequence converges also intrinsically in $D$.

Finally, a smooth enclosing cut in $D$ determines a filled competitor
containing the minimizing set $E$, hence also containing $O$. Its full
area is at least $P_g(E)=\Area_g(T)$ by global minimality. This is the
required outer area-minimizing property, with all components and any
frontier on $T$ retained.
\end{proof}

\subsection{Compactness of all minimizers and the area derivative}
\label{rig4:enc:variation-setup}

For a regular minimizing frontier $T=\partial E$, define its
unoriented tangent-plane area measure on the Grassmann bundle of
three-planes by
\begin{equation}
 \int\Phi\,dV_T
   =\int_T\Phi(x,T_xT)\,dA_g(x)
 \qquad
 \bigl(\Phi\text{ continuous on the Grassmann bundle}\bigr).
 \label{rig4:enc:tangent-measure}
\end{equation}
Here $V_T$ is a measure, not a volume form. Let $\mathcal T$ be the set
of all minimizing filled sets, identified almost everywhere and
recorded together with their regular frontiers. We give $\mathcal T$
the topology of local $L^1$ convergence of indicators and weak
convergence of the measures in Equation~\eqref{rig4:enc:tangent-measure}.
All variable parts and all these measures have support in one compact
region by Proposition~\ref{rig4:enc:minimizing-hull}.

\begin{proposition}[Compact minimizing space and right area derivative]
\label{rig4:enc:area-derivative}
The space $\mathcal T$ is compact and metrizable. If $h$ is a smooth
symmetric two-tensor, the function
\begin{equation}
 a'_T(h)=\frac12\int_T\tr_{T_xT}h\,dA_g
 \label{rig4:enc:metric-area-functional}
\end{equation}
is continuous on $\mathcal T$.

More generally, let $g_s$, $|s|<s_0$, be a smooth path of smooth metrics
on $X$, with $g_0=g$, uniformly comparable to $g$, and with uniformly
bounded first and second parameter derivatives relative to $g$.
Assume that their large coordinate spheres have strictly positive
outward mean curvature beyond one common radius. Write $a(s)$ for
their smooth-cut infimum and $h=\partial_sg_s|_{s=0}$. Then
\begin{equation}
 a'(0+)=\min_{T\in\mathcal T}a'_T(h).
 \label{rig4:enc:area-derivative-formula}
\end{equation}
In particular the minimum on the right is attained. If $s_j\to0$ and
$T_j$ is any minimizing frontier for $g_{s_j}$, a subsequence of its
filled sets converges to a member of $\mathcal T$, and its
$g$-tangent-plane area measures converge to the corresponding measure.
\end{proposition}

\begin{proof}
The proof of the outer barrier in
Proposition~\ref{rig4:enc:minimizing-hull} is uniform for this family:
all minimizing frontiers lie in one compact set $L\subset X$.
The metrics can be extended smoothly through the fixed boundary in a
common collar. On that collar and on $L$, their smooth dependence on
$s$ supplies uniform local bounds and convergence to the fixed extended
metric. The regularity and density results used above therefore remain
available with uniform local bounds where needed.

First consider a sequence of minimizers for $g$. BV compactness gives
$\mathbf 1_{E_j}\to\mathbf 1_E$ locally in $L^1$ along a subsequence,
with $E$ satisfying the same obstacle and exterior-tail constraints.
Lower semicontinuity and global minimality imply
\[
 a_g(B)\le P_g(E)\le\liminf_jP_g(E_j)=a_g(B).
\]
Thus $E$ is another minimizer, and the perimeters converge as well.

This strict BV convergence also gives the claimed tangent-plane
convergence, including mass carried by the original boundary. We give
the localization argument to keep track of the metric norm. Set
$\sigma_j=|D\mathbf 1_{E_j}|_g$ and
$\sigma=|D\mathbf 1_E|_g$, as scalar perimeter measures on
$\widetilde X$. Every subsequential weak limit of $\sigma_j$ dominates
$\sigma$: this follows from lower semicontinuity of perimeter weighted
by an arbitrary nonnegative continuous function. Their total masses
converge and their supports lie in one compact set, so this domination
must be equality. Hence $\sigma_j$ converges weakly to $\sigma$.
This conclusion uses only weighted lower semicontinuity, before any
continuity theorem for polar directions is applied.

Now form the vector-valued measures
$\mu_j=\nu_{E_j}\sigma_j$ and $\mu=\nu_E\sigma$.
Gauss--Green and local $L^1$ convergence give $\mu_j\rightharpoonup\mu$.
Choose finitely many smooth coordinate charts covering the common
compact support and a nonnegative smooth partition of unity
$\{\rho_\alpha\}$ subordinate to them. In each chart choose a smooth
matrix $A_\alpha(x)$ such that
$|A_\alpha(x)v|_\delta=|v|_g$. The coordinate vector measures
\[
 \lambda_{\alpha,j}=\rho_\alpha A_\alpha\mu_j
 \quad\hbox{satisfy}\quad
 |\lambda_{\alpha,j}|_\delta=\rho_\alpha\sigma_j.
\]
They converge weakly, and their Euclidean total variation masses
converge by the already established weak convergence of $\sigma_j$.
The vector-measure continuity theorem therefore applies in each chart
\cite[Theorem~1.3]{Spector2011}. Use its bounded continuous polar
integrand
\[
 (x,z)\longmapsto
 \Phi\bigl(x,(A_\alpha(x)^{-1}z)^{\perp_g}\bigr),
 \qquad |z|_\delta=1,
\]
and sum over the partition. This proves weak convergence of $V_{T_j}$
to $V_T$. Charts crossing $B$ are ordinary interior charts of
$\widetilde X$, so no contact mass is removed in this localization.
The compact base region and its Grassmann bundle are metrizable, so
the indicated $L^1$ and weak-measure topology is metrizable.
The subsequence argument proves compactness. Taking
$\Phi(x,\Pi)=\tfrac12\tr_\Pi h$ proves the asserted continuity.

For the varying metrics, uniform comparison with $g_s\to g$ on $L$
first gives
\begin{equation}
 a(s)\longrightarrow a(0),\qquad
 P_g(E_s)\longrightarrow a(0)
 \label{rig4:enc:varying-perimeters}
\end{equation}
for any choice of a minimizer $E_s$ for $g_s$. Indeed testing on a
fixed $g$-minimizer gives the upper bound for $a(s)$; comparison of
perimeters gives the lower bound, since $P_g(E_s)\ge a(0)$.
Consequently every BV limit of such $E_s$ is a $g$-minimizer and the
same strict-convergence argument gives tangent-plane convergence.
This proves the last assertion of the proposition.

It remains to calculate the derivative. The area element of a fixed
three-plane has the uniform Taylor expansion
\begin{equation}
 dA_{g_s}|_\Pi
 =\left(1+\frac{s}{2}\tr_\Pi h+O(s^2)\right)dA_g|_\Pi.
 \label{rig4:enc:area-taylor}
\end{equation}
The remainder is uniform over $x\in L$ and all three-planes $\Pi$,
because of the parameter-derivative bounds and uniform metric
positivity. Integrating over the fixed reduced boundary of a
competitor gives
\[
 P_{g_s}(E)=P_g(E)+s\,\frac12
       \int_{\partial^*E}\tr_{T_x\partial^*E}h\,dA_g
       +O(s^2)P_g(E).
\]
For each $T\in\mathcal T$, using its fixed filled set as a competitor
therefore yields
\[
 \limsup_{s\downarrow0}\frac{a(s)-a(0)}s
 \le a'_T(h).
\]
The continuous function in Equation~\eqref{rig4:enc:metric-area-functional}
has a minimum on $\mathcal T$, so this gives the required upper bound
by that minimum.

Conversely, choose a minimizer $E_s$ for each $s>0$. Its
$g$-perimeter is at least $a(0)$ and is uniformly bounded by
Equation~\eqref{rig4:enc:varying-perimeters}. Equation~\eqref{rig4:enc:area-taylor}
then gives
\[
 \frac{a(s)-a(0)}s
 \ge\frac12\int_{\partial E_s}
       \tr_{T_x\partial E_s}h\,dA_g-O(s).
\]
Along every sequence $s\downarrow0$, the compactness just proved gives
a subsequence whose integral converges to $a'_{T_\infty}(h)$ for some
$T_\infty\in\mathcal T$. This is at least the minimum in
Equation~\eqref{rig4:enc:area-derivative-formula}. The lower and upper
bounds agree, proving existence of the right derivative and the
formula. No uniqueness or smooth selection of the minimizing
frontier has been assumed.
\end{proof}

\begin{proposition}[A common tangent plane for tied minimizers]
\label{rig4:enc:common-plane}
Let
\[
 \mathcal L=\bigcup_{T\in\mathcal T}(T\setminus B)
 \subset\operatorname{int}X.
\]
If $x\in\mathcal L$ belongs to two minimizing frontiers, those frontiers have the
same tangent plane at $x$. Consequently $x\mapsto\Pi_x$, where
$\Pi_x=T_xT$ for any minimizing frontier through $x$, is a well-defined
continuous three-plane field on $\mathcal L$ with its relative
topology. In particular, for every smooth symmetric two-tensor $K$,
the function $x\mapsto\tr_{\Pi_x}K$ is single valued and continuous
there.
\end{proposition}

\begin{proof}
Let $E$ and $F$ be minimizing filled sets. Their union and intersection
are admissible, and perimeter submodularity gives
\[
 P_g(E\cup F)+P_g(E\cap F)
 \le P_g(E)+P_g(F)=2a_g(B).
\]
Each term on the left is at least $a_g(B)$, so both sets are also
minimizers. If their original frontiers met with different tangent
planes at a point outside $B$, the union and intersection would have
a corner there: in smooth local graph coordinates, their boundary is
the minimum or maximum of two functions with distinct first
derivatives. That contradicts the $C^1$ regularity of the minimizing
frontiers. Hence the tangent planes agree.

We also need continuity as the minimizing frontier varies. Suppose
$x_j\in T_j\setminus B$ and $x_j\to x\in\mathcal L$.
Proposition~\ref{rig4:enc:area-derivative} gives, after passage to a
subsequence, a minimizing limit $T_\infty$ with convergence of
indicators and tangent-plane area measures. Choose a small ball about
$x$ with closure disjoint from $B$. Uniform local perimeter density
bounds at $x_j$ imply that every smaller ball about $x$ carries
positive limiting perimeter, so $x\in T_\infty$.

On this ball every compactly supported perimeter replacement is
admissible, since the ball misses the obstacle. Thus each frontier is
an ordinary local area-minimizing boundary. Its associated integral
varifold is stationary for the fixed metric, has density one, and has
no boundary in the ball. The limiting frontier $T_\infty$ is smooth
and has multiplicity one.

These facts supply more than strict BV convergence alone. Choose a
sufficiently small continuity radius about $x$. Smoothness of
$T_\infty$ makes its normalized mass and its height and tilt excesses
as close as desired to those of the single plane $T_xT_\infty$.
Varifold convergence transfers these bounds to $T_j$ for large $j$;
moving the centers to $x_j$ changes the enclosing radii by a quantity
tending to zero. The rescaled smooth-metric errors also tend to zero
as the chosen radius decreases. Stationarity gives the required
mean-curvature bound in the Riemannian small-excess regularity theorem.
Equivalently, after a smooth local isometric embedding the Euclidean
generalized mean curvature is bounded by the fixed ambient second
fundamental form, whose rescaled contribution tends to zero.

The small-excess graphical regularity theorem consequently describes
the entire support in a smaller ball as one graph, with uniform
$C^{1,\alpha}$ bounds and tangent deviation tending to zero with
these inputs \cite{Allard1972,Maggi2012}; see \cite[Chapter~5, \S5, Theorem~5.2, p.~121]{Simon2014}. It cannot leave an additional
small sheet or spike outside that description. Oppositely cooriented
sheets also contribute positively to varifold mass and are covered
by the same conclusion. Express the graphs over the fixed plane
$T_xT_\infty$ on a still smaller disk. Their uniform $C^{1,\alpha}$
bounds give subsequential $C^{1,\beta}$ convergence for
$0<\beta<\alpha$, and varifold convergence identifies every such
graphical limit with the corresponding portion of $T_\infty$.
Thus the whole graph sequence converges in $C^1$ there, which in
particular gives
\[
 T_{x_j}T_j\longrightarrow T_xT_\infty.
\]
The already proved common-plane assertion identifies the plane on the
right with $\Pi_x$. Every convergent subsequence has this same plane
limit; compactness of the Grassmannian then proves continuity for the
whole sequence. The assertion about $\tr_{\Pi_x}K$ follows at once.
\end{proof}
